% !TeX root = ../thesis.tex
% ===================== ГЛАВА 2 =====================
\chapter{Верификация солвера на задачах релаксации}

Для проверки корректности и точности солвера рассмотрены две пространственно-однородные задачи релаксации, имеющие ясные
физические критерии: сохранение инвариантов, монотонность $H$-функции и установление равновесного распределения
Максвелла.

\section{Релаксация анизотропного распределения (квази-шок)}

Рассматривается релаксация пространственно-однородного неравновесного состояния, моделирующего столкновение двух
встречных потоков газа.
Начальное условие~--- бимодальное распределение
\begin{equation}
  f_0=\frac{0{,}5}{(2\pi)^{3/2}}\Big[
      e^{-\frac{(\xi_x-u)^2+\xi_y^2+\xi_z^2}{2}}
      +e^{-\frac{(\xi_x+u)^2+\xi_y^2+\xi_z^2}{2}}\Big],
\end{equation}
с тремя значениями параметра $u\in\{0{,}5;\,1{,}0;\,1{,}5\}$.
Сетка скоростей: $N=20$, $\xi_{\mathrm{cut}}=4{,}8$.
Требуется проследить эволюцию к равновесию и построить графики температуры $T(t)$, продольной компоненты тензора
температуры $T_{xx}(t)$ и $H$-функции.

Расчёт показал (рис.~\ref{fig:t0T}--\ref{fig:t0H}): общая температура $T$ сохраняется с машинной точностью (контроль
энергии), продольная компонента $T_{xx}$ релаксирует к $T$ за счёт «перекачки» энергии в поперечные степени свободы, а
$H$-функция монотонно убывает в соответствии с $H$-теоремой.
Доля пропущенных по условию положительности точек оказалась нулевой для всех трёх значений $u$, что с большим запасом
укладывается в допустимый порог.

\begin{figure}[htbp]
  \centering
  \includegraphics[width=0.7\textwidth]{t0_T}
  \caption{Эволюция температуры $T(t)$ для трёх значений $u$: общая температура сохраняется (контроль энергии).}
  \label{fig:t0T}
\end{figure}

\begin{figure}[htbp]
  \centering
  \includegraphics[width=0.7\textwidth]{t0_aniso}
  \caption{Затухание анизотропии $|T_{xx}-T|$ в логарифмическом масштабе в сравнении с моделью БГК.}
  \label{fig:t0aniso}
\end{figure}

\begin{figure}[htbp]
  \centering
  \includegraphics[width=0.7\textwidth]{t0_H}
  \caption{Эволюция $H$-функции Больцмана: монотонное невозрастание ($H$-теорема).}
  \label{fig:t0H}
\end{figure}

Для интерпретации результат сравнивался с модельным уравнением Бхатнагара--Гросса--Крука (БГК), которое в линейном
режиме даёт экспоненциальную релаксацию $f(t)=f_M+(f(0)-f_M)\,e^{-t/t_{\mathrm{rel}}}$ с единым временем релаксации.
Для модели твёрдых сфер в принятой нормировке при опорном состоянии ($n=1$, $T=1$) оно составляет~\cite{Tcher}
\begin{equation}
  t_{\mathrm{rel}}=\frac{5\sqrt{2}}{16}\approx0{,}442.
\end{equation}
Удобной мерой служит продольный момент четвёртого порядка $M_{xxxx}=\frac1n\sum_\alpha\xi_{x,\alpha}^4 f_\alpha\Delta
V$, равный в равновесии $3T^2$.
Как видно из рис.~\ref{fig:t0m4}, модель БГК с одним временем релаксации упрощает многомодовую структуру реального
оператора столкновений: моменты разных порядков релаксируют со своими скоростями, и наблюдаемое затухание анизотропии
оказывается в среднем медленнее экспоненты БГК, причём на ранних временах для больших $u$ кривая может вести себя
немонотонно.

\begin{figure}[htbp]
  \centering
  \includegraphics[width=\textwidth]{t0_m4}
\caption{Продольный $M_{xxxx}$ и поперечный $M_{yyyy}$ моменты 4-го порядка: многомодовость релаксации против модели БГК
(пунктир).}
  \label{fig:t0m4}
\end{figure}

Наконец, эволюция одномерной функции распределения $f_x(\xi_x)=\iint f\dd\xi_y\dd\xi_z$ (рис.~\ref{fig:t0fmarg})
наглядно показывает переход от двугорбого распределения (два пика на $\xi_x=\pm u$~--- две встречные «максвелловские
горки») к одногорбому симметричному распределению с максимумом в нуле, то есть к равновесному Максвеллу.
Усреднение по поперечным компонентам опирается на симметрию и подавляет численный шум.

\begin{figure}[htbp]
  \centering
  \includegraphics[width=0.95\textwidth]{t0_fmarg}
  \caption{Эволюция одномерной функции распределения $f_x(\xi_x)$ от двугорбой к равновесной максвелловской.}
  \label{fig:t0fmarg}
\end{figure}

Количественная сводка расчёта приведена в табл.~\ref{tab:t0}.
Для всех трёх значений $u$ общая температура $T$ совпадает с установившимся значением $T_{xx}$ (контроль энергии),
$H$-функция убывает, а доля пропущенных по положительности точек равна нулю~--- алгоритм устойчив при выбранных $\tau$ и
$p$.

\begin{table}[h]
  \centering
\caption{Сводка по релаксации квази-шока: число шагов до сходимости, продольная компонента $T_{xx}$ в начале и конце,
равновесная температура $T$ и значения $H$-функции.}
  \label{tab:t0}
  \begin{tabular}{cccccc}
    \hline
    $u$ & Шагов & $T_{xx}$ (старт) & $T_{xx}$ (конец) & $T$ & $H$: старт\,$\to$\,конец\\
    \hline
    0,5 & 419 & 1,249 & 1,083 & 1,083 & $-4{,}367\to-4{,}376$\\
    1,0 & 440 & 1,995 & 1,331 & 1,331 & $-4{,}591\to-4{,}685$\\
    1,5 & 421 & 3,225 & 1,740 & 1,740 & $-4{,}773\to-5{,}083$\\
    \hline
  \end{tabular}
\end{table}

\section{Релаксация бинарной смеси}

Вторая задача~--- релаксация пространственно-однородной смеси двух одноатомных газов с равными концентрациями
$n_c=n_h=0{,}5$ и максвелловскими начальными распределениями, температуры которых различаются вдвое: «холодный» газ
$T_c=1$ и «горячий» $T_h=2$.
Каждая компонента эволюционирует под действием интеграла столкновений со всей смесью $f=f_c+f_h$; требуется вычислить
общую и парциальные температуры, концентрации и энтропии.
Сетка скоростей выбрана крупнее, чтобы захватить тепловой хвост горячей компоненты: $N=28$, $\xi_{\mathrm{cut}}=7{,}0$.

Установившаяся температура смеси есть
\begin{equation}
  T_{\mathrm{eq}}=\frac{n_c T_c+n_h T_h}{n_c+n_h}=1{,}5,
\end{equation}
и расчёт воспроизводит это значение: парциальные температуры монотонно сближаются к общему значению
(рис.~\ref{fig:t1T}), концентрации компонент сохраняются с точностью округления $n_c=n_h=0{,}5$ (рис.~\ref{fig:t1n})~---
прямой контроль работы метода, а суммарная энтропия $S=-\sum f\ln f\,\Delta V$ монотонно возрастает
(рис.~\ref{fig:t1S}).
При этом холодный газ нагревается (его энтропия растёт), а горячий остывает (его энтропия убывает), тогда как энтропия
смеси в целом растёт, достигая максимума в равновесии~--- в согласии со вторым началом термодинамики.
Расчёт выполнен за 310 шагов по времени; доля пропущенных по положительности столкновений практически нулевая.

\begin{figure}[htbp]
  \centering
  \includegraphics[width=0.7\textwidth]{t1_T}
  \caption{Выравнивание парциальных температур $T_c$ и $T_h$ к общему значению $T_{\mathrm{eq}}=1{,}5$.}
  \label{fig:t1T}
\end{figure}

\begin{figure}[htbp]
  \centering
  \includegraphics[width=0.7\textwidth]{t1_n}
  \caption{Сохранение концентраций компонент $n_c=n_h=0{,}5$ (контроль метода).}
  \label{fig:t1n}
\end{figure}

\begin{figure}[htbp]
  \centering
  \includegraphics[width=0.7\textwidth]{t1_S}
  \caption{Энтропии: суммарная энтропия смеси монотонно возрастает ($H$-теорема); компоненты ведут себя противоположно.}
  \label{fig:t1S}
\end{figure}

Количественная сводка приведена в табл.~\ref{tab:t1}: парциальные температуры сходятся к общему значению $T=1{,}5$,
концентрации компонент сохраняются точно, а суммарная энтропия монотонно растёт.

\begin{table}[h]
  \centering
\caption{Сводка по релаксации смеси: парциальные и общая температуры, концентрации и суммарная энтропия в начальный и
конечный (равновесный) моменты.}
  \label{tab:t1}
  \begin{tabular}{lcc}
    \hline
    Величина & $t=0$ & равновесие\\
    \hline
    $T_c$ (холодный газ) & 1,000 & 1,500\\
    $T_h$ (горячий газ)  & 2,000 & 1,500\\
    $T$ (смесь)          & 1,500 & 1,500\\
    $n_c=n_h$            & 0,500 & 0,500\\
    $S$ (смесь)          & 4,855 & 4,865\\
    \hline
  \end{tabular}
\end{table}

Совокупность результатов Главы~2~--- точное сохранение инвариантов, монотонность $H$-функции, установление равновесного
распределения Максвелла и правильное равновесие смеси~--- подтверждает корректность реализованного солвера и позволяет
применить его к неоднородной задаче переноса.

